\documentclass[12pt]{article}
\usepackage[T1]{fontenc}
\usepackage[utf8]{inputenc}
\usepackage{lmodern}
\usepackage{textcomp}
\usepackage{enumitem}
\usepackage{geometry}
\geometry{margin=1in}
\begin{document}
\begin{center}
Answer Key - History - K-12 Level Assessment 19 \
\small (Answers are ordered exactly as in the question paper.)
\end{center}

\section*{Multiple Choice}
\begin{enumerate}[label=\arabic*.]
\item \textbf{Answer:} C
\\ \textit{Options:} A) sunflower; B) squash; C) maize; D) knotweed
\\ \textit{Note:} Correct option: C - maize
\item \textbf{Answer:} D
\\ \textit{Options:} A) many processes that can be observed act extremely slowly.; B) uniformitarianism can account for significant geological features.; C) erosion works incredibly slowly.; D) All of the above.
\\ \textit{Note:} Correct option: D - All of the above.
\item \textbf{Answer:} C
\\ \textit{Options:} A) his horses and guns frightened the Incas into submission.; B) he was such a brilliant military strategist.; C) the empire had never won the allegiance of the peoples it had conquered.; D) he arrived with a superior force of more than 80,000 soldiers.
\\ \textit{Note:} Correct option: C - the empire had never won the allegiance of the peoples it had conquered.
\item \textbf{Answer:} B
\\ \textit{Options:} A) Khipu; B) Mit'a; C) Split inheritance; D) Capacocha
\\ \textit{Note:} Correct option: B - Mit'a
\item \textbf{Answer:} C
\\ \textit{Options:} A) the use of seeds as a natural insecticide in Altamira cave in southwestern Europe; B) plant materials such as flowers in Aurignacian burials in southern France; C) plant material lining the floor of Sibudu Cave in South Africa; D) the use of wood in the Gravettian tradition to make spear-throwers and canoes
\\ \textit{Note:} Correct option: C - plant material lining the floor of Sibudu Cave in South Africa
\end{enumerate}

\section*{True / False}
\begin{enumerate}[label=\arabic*.]
\setcounter{enumi}{5}
\item \textbf{Answer:} False
\\ \textit{Options:} A) True; B) False
\\ \textit{Note:} LLM-generated false statement. Correct answer: 'Clovis'.
\item \textbf{Answer:} True
\\ \textit{Options:} A) True; B) False
\\ \textit{Note:} LLM-generated true statement based on MCQ.
\item \textbf{Answer:} False
\\ \textit{Options:} A) True; B) False
\\ \textit{Note:} LLM-generated false statement. Correct answer: '20,000 B.P.'.
\item \textbf{Answer:} False
\\ \textit{Options:} A) True; B) False
\\ \textit{Note:} LLM-generated false statement. Correct answer: 'Savagery, Barbarism, and Civilization.'.
\item \textbf{Answer:} False
\\ \textit{Options:} A) True; B) False
\\ \textit{Note:} LLM-generated false statement. Correct answer: 'uniformitarianism.'.
\end{enumerate}

\section*{Long Form Response}
\begin{enumerate}[label=\arabic*.]
\setcounter{enumi}{10}
\item \textbf{Answer:} Paleoindians inhabited a wide diversity of habitats, including some of the highest elevations in the New World.. Explain why this is the correct answer and provide supporting evidence. Reference key facts or concepts that support this choice.
\\ \textit{Note:} Long-form derived from MCQ; award credit for stating the correct idea and giving supporting reasoning.
\item \textbf{Answer:} They brought animals to the site from outside of the region.. Explain why this is the correct answer and provide supporting evidence. Reference key facts or concepts that support this choice.
\\ \textit{Note:} Long-form derived from MCQ; award credit for stating the correct idea and giving supporting reasoning.
\end{enumerate}

\vfill
\noindent\textit{End of Answer Key}
\end{document}